\begin{lemma}[Universal calibration and fixed log-star threshold]\label{lem:step-005-calibration}
Under the universal constants \(b_*,d_*>0\) and
\(N_*\in\mathbb Z_{\ge2}\) supplied by Proposition~\ref{prop:step-002-wrapper}, define

\[
q_*:=\left\lceil\frac{16}{b_*}\right\rceil,
\qquad
N_0:=\max\left\{N_*,\operatorname{Tow}_2(q_*)\right\},
\tag{1}
\]

and choose the constants in \eqref{eq:calibration}. Then

\[
a,c_\delta,\varepsilon_0>0,\qquad
\alpha_0,\beta_0\in(0,1/2),qquad
N_0\in\mathbb Z_{\ge2},
\]

all choices are universal and independent of \(k,N,n,\varepsilon,\delta\),
and every \(N\ge N_0\) satisfies

\[
N\ge N_*,
\qquad
b_*\log_2^*N\ge16.
\tag{2}
\]

In particular, (2) holds at the boundary \(N=N_0\).
\end{lemma}

\begin{proof}
Because \(b_*>0\), \(q_*\) is an integer at least one and

\[
q_*\ge\frac{16}{b_*}.
\tag{3}
\]

The tower convention gives

\[
\log_2^*\operatorname{Tow}_2(q)=q
\qquad(q\in\mathbb Z_{\ge1}).
\tag{4}
\]

Indeed, the claim is immediate for \(q=1\), since
\(\operatorname{Tow}_2(1)=2\). If it holds at \(q\), then the first
base-two logarithm of
\(\operatorname{Tow}_2(q+1)=2^{\operatorname{Tow}_2(q)}\) is
\(\operatorname{Tow}_2(q)\), so exactly one further iteration is needed.

More directly for the comparison needed here, if
\(N\ge\operatorname{Tow}_2(q_*)\), then for every
\(0\le j<q_*\), monotonicity of the iterated base-two logarithm on the
positive range traversed here gives

\[
\log_2^{(j)}N
\ge
\log_2^{(j)}\operatorname{Tow}_2(q_*)
>1.
\]

Hence \(L_N=\log_2^*N\ge q_*\). Combining this with (3) yields

\[
b_*L_N\ge b_*q_*\ge16.
\tag{5}
\]

Definition (1) also gives \(N\ge N_*\), including when \(N=N_0\).
The remaining assertions follow from the positivity of
\(b_*,d_*\), the exact choices
\(a=b_*/16\), \(c_\delta=d_*\), \(\varepsilon_0=0.1\), and
\(2^{-13}\in(0,1/2)\). Since \(b_*,d_*,N_*\) are universal by
Proposition~\ref{prop:step-002-wrapper},
so are \(q_*\), \(N_0\), and every constant in
\eqref{eq:calibration}. \end{proof}

\begin{lemma}[Strict domination of the exact simulated budget]\label{lem:step-005-budget}
Under Assumption~\ref{assump:candidate-regime} and
Lemma~\ref{lem:step-005-calibration}, if

\[
n<akL_N
\qquad\text{with}\qquad a=\frac{b_*}{16},
\]

then the exact setting-defined budget

\[
M=\max\left\{8,\left\lceil\frac{4n}{k}\right\rceil\right\}
\]

is an integer satisfying

\[
8\le M<b_*L_N.
\tag{6}
\]

This remains strict when the maximum is attained by the floor-eight
branch, when \(4n/k\) is an integer, and at \(N=N_0\).
\end{lemma}

\begin{proof}
The contradiction hypothesis and \(k>0\) give the strict inequality

\[
\frac{4n}{k}
<4aL_N
=\frac14 b_*L_N.
\tag{7}
\]

For every real \(x\), including integral \(x\),
\(\lceil x\rceil<x+1\). By (5), \(1\le b_*L_N/16\). Therefore

\[
\begin{aligned}
\left\lceil\frac{4n}{k}\right\rceil
&<\frac{4n}{k}+1\\
&<\frac14b_*L_N+1\\
&\le\frac{5}{16}b_*L_N\\
&<b_*L_N.
\end{aligned}
\tag{8}
\]

The fixed lower branch is also strictly below the source threshold:

\[
8<16\le b_*L_N.
\tag{9}
\]

Both entries of the maximum are strictly below \(b_*L_N\), so their
maximum is strictly below it. Since both entries are integers and the first
is 8, \(M\in\mathbb Z_{\ge8}\), proving (6).

At the requested smallest tuple \((n,k)=(1,2)\),

\[
\left\lceil\frac{4n}{k}\right\rceil=\lceil2\rceil=2,
\qquad M=8.
\]

Whenever the local contradiction hypothesis holds, (9) proves the strict
hard-regime inequality for this tuple as well. At \(N=N_0\), (5) still
holds, so no asymptotic or strictly-larger-domain exception is being used.
\end{proof}

\begin{lemma}[Candidate-conjunction preservation and strict source privacy cap]\label{lem:step-005-privacy-cap}
Under Assumption~\ref{assump:candidate-regime}, with
\(\varepsilon_0=0.1\), \(0<c_\delta\le d_*\), and the exact
\(M=m_{n,k}\ge8\), one has simultaneously

\[
0<\varepsilon\le0.1,
\tag{10}
\]

\[
0<\delta\le\frac{1}{n\log(n+1)},
\tag{11}
\]

\[
0<\delta\le\frac{c_\delta}{M^2\log(M+1)}
<\frac{d_*}{M^2\log M}.
\tag{12}
\]

Thus both original \(\delta\)-conjuncts are retained, while the second
strictly implies the source cap.
\end{lemma}

\begin{proof}
Equation (10) is the candidate \(\varepsilon\)-condition with the fixed
choice \(\varepsilon_0=0.1\). Since an upper bound by a minimum is an upper
bound by each entry, the candidate \(\delta\)-condition gives (11) and the
first inequality in (12) separately. Their denominators are well-defined
and positive: \(n\ge1\) gives \(\log(n+1)\ge\log2>0\), while \(M\ge8\)
gives

\[
\log(M+1)>\log M\ge\log8>0.
\tag{13}
\]

The first inequality in (13) is strict because the natural logarithm is
strictly increasing and \(M+1>M\). Taking positive reciprocals reverses
that strict comparison. Using \(0<c_\delta\le d_*\),

\[
\frac{c_\delta}{M^2\log(M+1)}
<\frac{c_\delta}{M^2\log M}
\le\frac{d_*}{M^2\log M},
\tag{14}
\]

which proves the strict part of (12). Notice that (11) was not used to
derive (14) and has not been replaced by it.

At the exact floor boundary \(M=8\), (12) reads

\[
0<\delta\le\frac{c_\delta}{64\log9}
<\frac{d_*}{64\log8}.
\tag{15}
\]

At \((n,k)=(1,2)\), where Lemma~\ref{lem:step-005-budget} computes
\(M=8\), the full primitive conjunction is explicitly

\[
0<\delta\le\frac1{\log2}
\quad\text{and}\quad
0<\delta\le\frac{c_\delta}{64\log9},
\]

and (15) supplies the source cap. \end{proof}

\begin{proposition}[Hard-regime and candidate-parameter certificate]\label{prop:step-005-certificate}
Under Assumption~\ref{assump:candidate-regime}, Proposition~\ref{prop:step-002-wrapper}, and the local contradiction
hypothesis \(n<ak\log_2^*N\), the universal choices in
\eqref{eq:calibration} and (1) imply
the following complete certificate.  The exact \(M=m_{n,k}\) satisfies

\[
N\ge N_*,
\qquad 8\le M<b_*\log_2^*N,
\qquad 0<\varepsilon\le0.1,
\]

\[
0<\delta\le\frac1{n\log(n+1)},
\qquad
0<\delta\le\frac{c_\delta}{M^2\log(M+1)}
<\frac{d_*}{M^2\log M}.
\]

The conclusion includes \(N=N_0\), \(M=8\), and, when the contradiction
hypothesis is active, \((n,k)=(1,2)\).
\end{proposition}

\begin{proof}
Lemma~\ref{lem:step-005-calibration} fixes the universal theorem constants
once and for all and gives \(N\ge N_*\) and \(b_*L_N\ge16\), including at
\(N=N_0\). Lemma~\ref{lem:step-005-budget} applies the local contradiction
hypothesis to the exact ceiling-and-maximum definition of \(M\), proving
\(M\in\mathbb Z_{\ge8}\) and the strict inequality \(M<b_*L_N\).
Lemma~\ref{lem:step-005-privacy-cap} gives \(0<\varepsilon\le0.1\), keeps
both candidate-$\delta$ inequalities, and proves the stronger strict
source-cap membership \(0<\delta<d_*/(M^2\log M)\).
The explicit boundary computations (9) and (15), together with the
\((n,k)=(1,2)\) calculation after (9), prove the stated boundary clauses.
\end{proof}
